\centerline{\bf \Huge Клетачки)))0)0))}

\begin{flushright}
        \scriptsize{}
\end{flushright}

\problem{1}
$\bigl[\textbf{Лига Профессионалов, 2024}\bigr]$ 
Даны нечетные натуральные $m, n>3$. Клетки доски $m \times n$ окрашены в чёрный и белый цвет. Обозначим через $A$ количество строк, в которых чёрных клеток больше, чем белых, через $B$ - количество столбцов, в которых белых клеток больше, чем чёрных. Найдите наибольшее возможное значение $A+B$.

\problem{2}
$\bigl[\textbf{УТЮМ, 2023}\bigr]$
Каждая клетка квадрата $n \times n$ может быть покрашена в белый или черный цвет. За один ход можно выбрать 4 клетки, являющиеся вершинами прямоугольника со сторонами, параллельными линиям сетки, и перекрасить их в противоположный цвет. Для какого наибольшего $k$ можно найти $k$ различных раскрасок квадрата так, чтобы ни одну из выбранных раскрасок нельзя было бы привести ни к одной другой с помощью выбранных операций?

\problem{3}
$\bigl[\textbf{ММО, 2020}\bigr]$
Для каких $k$ можно закрасить на белой клетчатой плоскости несколько клеток (конечное число, большее нуля) в черный цвет так, чтобы на любой клетчатой вертикали, горизонтали и диагонали либо было ровно $k$ черных клеток, либо вовсе не было черных клеток?

\problem{4}
$\bigl[\textbf{Закл, 2016}\bigr]$
Из клетчатого бумажного квадрата $100 \times 100$ вырезали по границам клеток 1950 двуклеточных прямоугольников. Докажите, что из оставшейся части можно вырезать по границам клеток четырехклеточную фигурку в виде буквы T , возможно, повернутую. (Если такая фигурка уже есть среди оставшихся частей, считается, что ее получилось вырезать.)

\problem{5}
$\bigl[\textbf{Поволжская олимпиада, 2020}\bigr]$ 
В какое наибольшее количество цветов можно покрасить клетки квадрата $50 \times 50$ так, чтобы в любом квадрате $2 \times 2$ были хотя бы две одноцветные клетки?


\newpage

\hspace{3mm}